\vfill\pagebreak
\section{Putting it together: a calendar}
\label{sec:functions:calendar}

The program of this section prints the calendar of a month, or of a whole year,
laid out week by week like the calendar on a wall. It is the first program of the
book long enough to need functions, and it is a useful one: the calendar of any
month from 1900 on, in a terminal.

\subsection{Splitting the problem}

A program of this size is not written from the first line to the last. It is
easier to start from what it must do, and to split that into questions small
enough to be answered by one short function each. Printing the calendar of a
month raises these questions:

\begin{itemize}
\item \textit{What is the month called?} A \token{match} on the month number
  answers it.
\item \textit{How many days does the month have?} Exercise~6 of
  Chapter~\ref{chap:control} answered it for ordinary years. February has 29 days
  in a leap year, which raises the next question.
\item \textit{Is the year a leap year?} It is when it is divisible by 4, except
  for the years divisible by 100, which are leap years only when they are also
  divisible by 400: 2024 and 2000 are leap years, 1900 is not.
\item \textit{On which day of the week does the month start?} This decides how
  far to the right the 1st is printed. The program knows that January 1st, 1900
  was a Monday. It counts the days between that date and the 1st of the month,
  using the number of days in each year and in each month on the way, and the
  remainder of their division by 7 is the day of the week.
\end{itemize}

Each question becomes a function, and the functions that answer the last
questions call those that answer the first ones. Printing the calendar then
takes only a few lines, because the hard parts have names.

\subsection{The program}

\lstinputlisting[style=coloredverbatim, caption={A calendar, \textit{examples/functions/calendar.yr}}, label=lst:functions:calendar]{examples/functions/calendar.yr}
\vspace{-10pt}%
\begin{lstlisting}[style=bashVerb, escapechar=@+]
@+\textcolor{teal!80}{alice@dev:\mtilde}@+$ gyc calendar.yr -o calendar
@+\textcolor{teal!80}{alice@dev:\mtilde}@+$ ./calendar
Year: 1850
Please type a number from 1900 to 9999.
Year: 2026
Month (0 for the whole year): 9

     September 2026
Mo Tu We Th Fr Sa Su
    1  2  3  4  5  6
 7  8  9 10 11 12 13
14 15 16 17 18 19 20
21 22 23 24 25 26 27
28 29 30
\end{lstlisting}

\subsection{How it works}

The functions are declared in the order of the questions, each after the ones it
calls. Nothing requires that order, but it lets the file be read from top to
bottom.

\begin{itemize}
\item \token{isLeapYear}, on line~3, is one Boolean expression, whose value is
  the answer.
\item \token{daysIn} is overloaded. With one argument, on line~7, it gives the
  number of days in a year; with two, on line~11, the number of days in a month
  of a given year. Both are ``the number of days in'' something, and the
  arguments say what. The arm for February is an \token{if} used as a value.
\item \token{weekday}, on line~38, counts the days from January 1st, 1900. The
  first loop adds the days of every year before \token{year}, the second the days
  of every month before \token{month}, and line~46 adds the days of the month
  before \token{day}. Its three parameters have the same type and a day, a month
  and a year are written in different orders in different countries, so all
  three are marked \token{@}. The call on line~59 cannot mix them up.
\item \token{printDay}, on line~49, prints a day on two characters, padded with a
  space, so that the columns line up.
\item \token{printCalendar}, on line~56, prints the title and the header, then
  three spaces for each day of the week before the 1st, then the days. After a
  day, the number of cells printed on the current row is \token{first + day},
  counting the blank ones; when it is a multiple of 7, the week is complete and
  the row ends. The \token{if} on line~70 ends the last row, unless the month
  ended on a Sunday and the row is already ended.
\item \token{printCalendar} is overloaded too. The version with one parameter,
  on line~75, prints a whole year by calling the other for each month.
\item \token{askNumber}, on line~82, asks the same question until the answer is
  within bounds. Its bounds must be named, so the call on line~93 reads as
  ``a number from 1900 to 9999''. The \token{loop} produces the number accepted,
  with \token{break n} (Section~\ref{sec:control:loop}).
\item \token{main} is left with nothing to do but ask and choose.
\end{itemize}

Every function of the program is short enough to be checked on its own. If a
month started on the wrong day, only \token{weekday} could be at fault, and it
can be tested alone, with a few dates whose day is known. The last exercise of
the chapter reuses it.
